\documentclass[10pt,a4paper]{amsart}
\usepackage[margin=19mm]{geometry}
\usepackage{amsmath,amssymb,mathtools}
\usepackage[scr=rsfs,cal=euler]{mathalfa}
\usepackage{xfrac}
\usepackage[unicode,hidelinks,bookmarks=false]{hyperref}
\newcommand{\ie}{i.e.\ }
\newcommand{\cf}{cf.\ }
\newcommand{\resp}{respectively\ }
\newcommand{\Subsection}{\subsection}
\providecommand{\nobreakdash}{\nobreak\hbox{-}}
\setlength{\emergencystretch}{2em}
\multlinegap0pt
\usepackage{bm}
\usepackage{bbm}
\usepackage{dsfont}
\usepackage{xspace}
\usepackage{cancel}
\usepackage{mathtools}
\usepackage{amsthm}
\makeatletter
\@ifundefined{theorem}{\newtheorem{theorem}{Theorem}}{}
\@ifundefined{lemma}{\newtheorem{lemma}[theorem]{Lemma}}{}
\@ifundefined{proposition}{\newtheorem{proposition}[theorem]{Proposition}}{}
\makeatother
\title[A variation of the inverse Fueter theorem and the generalize]{A variation of the inverse Fueter theorem and the generalized polyanalytic Cauchy-Kovalevskaya extension of order 2}
\author{Antonino De Martino, Stefano Pinton}
\date{Source: arXiv:2608.07711v1 (CC BY 4.0)}
\begin{document}
\maketitle

\noindent\emph{Source and licence.} A variation of the inverse Fueter theorem and the generalized polyanalytic Cauchy-Kovalevskaya extension of order 2. Version arXiv:2608.07711v1 (CC BY 4.0), distributed under the Creative Commons Attribution 4.0 International licence, as recorded in the arXiv licence metadata for this exact version and shown on the abstract page https://arxiv.org/abs/2608.07711v1. Full rights notice: licenses/arxiv-2608.07711-CC-BY-4.0.txt.\par\smallskip

\noindent\textit{Editorial scope.} Counted unit: the Proposition whose source label is \texttt{decop} in the article (source0.tex lines 1081--1093, source label \texttt{decop}), delivered together with the article's own complete original proof of it (source0.tex lines 1094--1188, one \texttt{proof} environment of 95 lines). No mathematics is written, repaired or re-derived: statement and proof are the article's own text line for line. Required context retained uncounted: decof (lines 561--574); seriesexp (lines 825--828); Stefano1 (lines 835--838); f1 (lines 888--898); pgck (lines 932--967); fun1 (lines 935--938); funct (lines 940--943); real (lines 945--950); imm (lines 952--957). Every definition, display and label of those lines is retained unchanged. Omitted from this package and not redistributed: 1--560, 575--824, 829--834, 839--887, 899--931, 939--939, 944--944, 951--951, 958--1080, 1189--2717. Nothing inside a retained excerpt is omitted or altered: every retained body is delivered exactly as the article's lines are written, so no omission receipt of any kind accompanies this package. Citations: the retained excerpts carry no citation command at all, so no bibliography entry of the article is transcribed and no reference is redistributed. Reference adaptation: none was needed and none was made; every reference inside the delivered unit resolves inside it, so no unresolved reference or citation is produced. Printed slips kept literal: no printed word, symbol, hypothesis or number of the counted statement or of its proof was altered. Layout and class: the article's own class is \texttt{amsart} with packages including xcolor, graphicx, hyperref, amscd, amsmath, amsthm, amssymb, latexsym, array, amsfonts, shadow, amsbsy, dsfont, mathtools; the wrapper is the lane's pinned \texttt{amsart} editorial wrapper at 10pt on A4, which restates the theorem-like environments the retained text needs and adds the article's own macro definitions that the retained text needs (none), with extra wrapper packages (bm, bbm, dsfont, xspace, cancel, mathtools). The delivered numbering is the unit's own. Assets: no retained span contains a figure environment, a graphics-inclusion command, a PDF-page inclusion, a table or a TikZ picture; no image file is copied into this package, the package declares no asset dependency and no image file is redistributed with it. Licence: version arXiv:2608.07711v1 of this article is distributed under the Creative Commons Attribution 4.0 International licence, as recorded in the arXiv licence metadata for this exact version and shown on the abstract page https://arxiv.org/abs/2608.07711v1. A full rights notice is delivered at licenses/arxiv-2608.07711-CC-BY-4.0.txt. Characters: every retained excerpt is pure ASCII, so the delivered engine, XeLaTeX with its default Latin Modern text font, reports no missing character.
\begin{proposition}
\label{decof}
Let $U \subseteq \mathbb{H}$ be an open set. A function $f \colon U \to \mathbb{H}$ is left (resp.\ right) axially polyanalytic Fueter regular of order $2$ if and only if it admits the representation
$$
f(q)= g_1(q)+ q_0\, g_2(q),
$$
where
\begin{equation}
	\label{g12}
	g_1(q)= f(q)- q_0\, Df(q),
	\qquad
	g_2(q)= Df(q).
\end{equation}
\end{proposition}

\begin{equation}
	\label{seriesexp}
	f(q)= \frac{\sqrt{\pi}}{2}\left( \sum_{j=0}^\infty \frac{(-1)^j |\underline{q}|^{2j} \partial_{q_0}^{2j} A_0(q_0)}{2^{2j}j! \Gamma \left(\frac{3}{2}+j\right)}+ \frac{\underline{q} }{2} \sum_{j=0}^\infty \frac{(-1)^j |\underline{q}|^{2j} \partial_{q_0}^{2j+1}A_0(q_0)}{2^{2j}j! \Gamma \left(\frac{5}{2}+j\right)} \right).
\end{equation}

\begin{equation}
\label{Stefano1}
\Omega_1=  \{(q_0, \underline{q}) \in \mathbb{R}^4 \, \, |\, \, q_0 \in \tilde{D}, \, \, \, |\underline{q}|< R(q_0)\},
\end{equation}

\begin{lemma}
	Let $ \ell  \in \mathbb{N}$, then for $q \in \mathbb{H}$ we have
	\begin{equation}
		\label{f1}
		\partial_{\underline{q}}(\underline{q}^\ell)= c(\ell) \underline{q}^{\ell-1}, \qquad c(\ell)=\begin{cases}
			- \ell, \quad \ell \, \, \hbox{even}\\
			-(\ell+2), \quad \ell \, \, \hbox{odd},
		\end{cases}
	\end{equation}
where $\partial_{\underline{q}}= e_1 \partial_{q_1}+e_2 \partial_{q_2}+e_2 \partial_{q_3}$.
\end{lemma}

\begin{theorem}
\label{pgck}
Let $D$ be a convex intrinsic complex domain. We set $\tilde{D}:=D \cap \mathbb{R}$. We consider two analytic functions $A_0(q_0)$, $A_1(q_0) \in \mathcal{A}(\tilde{D}) \otimes \mathbb{H}$. Then there exist a unique sequence of functions $ \{A_j\}_{j \in \mathbb{N}_0} \subset \mathcal{A}(\tilde{D}) \otimes \mathbb{H}$ such that the series
\begin{equation}
\label{fun1}
f(q)= \sum_{j=0}^{\infty} \underline{q}^jA_j(q_0)
\end{equation}
converges in axially symmetric $4$-dimensional neighbourhood $\Omega_1 \subset \mathbb{H}$ of $\tilde{D}$ and its sum is polyanalytic of order 2 (i.e. $D^2 f(q)=0$) in $\Omega_1$. Furthermore, the sum of $f$ is given by the axial expression
\begin{equation}
\label{funct}
f(q)= \alpha(q)+\underline{\omega} \beta(q), \quad \underline{\omega}=\frac{\underline{q}}{|\underline{q}|}
\end{equation}
where
\begin{eqnarray}
	\nonumber
	\alpha(q)&:=& \frac{\sqrt{\pi}}{2}\left(-\frac{3}{2}  \sum_{j=0}^{\infty} \frac{(-1)^j |\underline{q}|^{2j+2} \partial_{q_0}^{2j+1} A_1(q_0)}{2^{2 j} j! \Gamma \left(j+\frac{5}{2}\right)} + \frac{1}{2}  \sum_{j=0}^{\infty} \frac{(-1)^j |\underline{q}|^{2j+2} \partial_{q_0}^{2j+2}A_0(q_0)}{2^{2j} j! \Gamma \left(j+\frac{5}{2}\right)} \right.\\
	\label{real}
	&&\left.+ \sum_{j=0}^\infty \frac{(-1)^j |\underline{q}|^{2j} \partial_{q_0}^{2j} A_0(q_0)}{2^{2j}j! \Gamma \left(j+\frac{3}{2}\right)} \right),
\end{eqnarray}
and
\begin{eqnarray}
	\nonumber
	\beta(q)&:=& -  \frac{\sqrt{\pi}}{4}\left( \frac{3}{2} \sum_{j=0}^\infty \frac{(-1)^j |\underline{q}|^{2 j+3} \partial_{q_0}^{2 j+2}A_1(q_0)}{2^{2j}j! \Gamma \left(j+\frac{7}{2}\right)}- 3 \sum_{j=0}^\infty \frac{(-1)^j |\underline{q}|^{2j+1} \partial_{q_0}^{2j}A_1(q_0)}{2^{2j}j! \Gamma \left(j+\frac{5}{2}\right)} \right.\\
	\label{imm}
	&&\left. - \frac{1}{2} \sum_{j=0}^\infty \frac{(-1)^j |\underline{q}|^{2j+3} \partial_{q_0}^{2j+3} A_0(q_0)}{2^{2j} j! \Gamma \left(j+\frac{7}{2}\right)} \right).
\end{eqnarray}
The functions $A_0$ and $A_1$ are determined by the relations
\begin{equation}
\label{infun}
A_0(q_0)= \lim_{|\underline{q}| \to 0}f(q), \qquad
A_1(q_0)= -\frac{1}{3} \lim_{|\underline{q}| \to 0}\partial_{\underline{q}} f(q).
\end{equation}
The function in \eqref{funct} is called the GCK-extension of polyanalytic functions of order 2 of the couple $(A_0, A_1)$, and it is denoted by $PGCK[A_0,A_1](q)$. This extension operator defines an isomorphism between right modules:
$$ PGCK: (\mathcal{A}(\tilde{D}) \otimes \mathbb{H})^2 \to \mathcal{APA}_2^L(\Omega_1),$$
where the set $\Omega_1$ is defined as in \eqref{Stefano1}.
\end{theorem}	

\begin{proposition}
\label{decop}
	Let $f(q) = \sum_{j=0}^{\infty} \underline{q}^{\,j} A_j(q_0)$ be a function satisfying the same assumptions as in Theorem~\ref{pgck}. Then the GCK extension for polyanalytic functions of order~2 (see \eqref{funct}) can be expressed as
\begin{equation}
\label{Pgck}
f(q) = \mathrm{GCK}[g_1(q_0)](q) + q_0\, \mathrm{GCK}[g_2(q_0)](q),
\end{equation}
where
$$
g_1(q) = f(q) - q_0\, Df(q), \qquad g_2(q) = Df(q),
$$
where $\mathrm{GCK}$ denotes the GCK extension for axially Fueter regular functions, as defined in \eqref{seriesexp}.
\end{proposition}

\begin{proof}
By hypothesis we know that the function $f$ can be written as in \eqref{fun1}, thus by \eqref{f1} we can write the function $g_1(q)$ as
\begin{eqnarray*}
g_1(q)&=& \sum_{j=0}^\infty \underline{q}^j A_j(q_0)- q_0 D \left(\sum_{j=0}^\infty  \underline{q}^j A_j(q_0)\right)\\
&=&\sum_{j=0}^\infty \underline{q}^j A_j(q_0)- q_0 \left( \sum_{j=0}^\infty \underline{q}^j \partial_{q_0} A_j(q_0)+ \sum_{j=1}^\infty c(j) \underline{q}^{j-1} A_j(q_0)\right),
\end{eqnarray*}
Similarly, for the function $g_2(q)$ we have
$$ g_2(q)=\sum_{j=0}^\infty \underline{q}^j \partial_{q_0} A_j(q_0)+ \sum_{j=1}^\infty c(j) \underline{q}^{j-1} A_j(q_0),$$
where where the constant $c(j)$ has been defined in \eqref{f1}.
Thus the restrictions of the functions $g_1(q)$ and $g_2(q)$ to the real line are given by
\begin{equation}
\label{res1}
g_1(q_0)= \lim_{|\underline{q}| \to 0}g_1(q) =A_0(q_0)-q_0 \partial_{q_0} A_0(q_0)+3q_0 A_1(q_0),
\end{equation}
\begin{equation}
\label{res2}
g_2(q_0)=\lim_{|\underline{q}| \to 0} g_2(q)=\partial_{q_0} A_0(q_0)-3A_1(q_0).
\end{equation}
By Proposition \ref{decof} we know that the functions $g_1(q)$ and $g_2(q)$ are axially Fueter regular, so we can compute their GCK extensions for Fueter regular functions. We start from the scalar part of the GCK extension of the functions $g_1(q)$ and $g_2(q)$, that we denote by $GCK_0$. By using the Leibniz rule, for $j \in \mathbb{N}$, we obtain

\begin{eqnarray}
\nonumber
(|\underline{q}| \partial_{q_0})^{2 j} g_1(q_0)&=& (|\underline{q}| \partial_{q_0})^{2 j}A_0(q_0)-q_0 | \underline{q}|^{2 j} \partial_{q_0}^{2 j +1} A_0(q_0)-2 j | \underline{q}|^{2 j} \partial_{q_0}^{2 j} A_0(q_0)\\
\label{rev1}
&&+3 q_0| \underline{q}|^{2 j} \partial_{q_0}^{2 j} A_1(q_0)+6 j | \underline{q}|^{2 j} \partial_{q_0}^{2 j -1}A_1(q_0) .
\end{eqnarray}
By using the scalar part of \eqref{seriesexp} and performing a change of index from $j-1$ to $j$ in the third and last summations after the last equality, we obtain
\begin{eqnarray}
	\nonumber
GCK_0[g_1(q_0)](q)&=& \frac{\sqrt{\pi}}{2} \sum_{j=0}^\infty \frac{(-1)^j (|\underline{q}| \partial_{q_0})^{2j}g_1(q_0)}{2^{2j}j! \Gamma \left(\frac{3}{2}+j\right)}\\
\nonumber
&=& \frac{\sqrt{\pi}}{2} \sum_{j=0}^\infty \frac{(-1)^j (|\underline{q}| \partial_{q_0})^{2j}A_0(q_0)}{2^{2j}j! \Gamma \left(\frac{3}{2}+j\right)}-\frac{q_0\sqrt{\pi}}{2} \sum_{j=0}^\infty \frac{(-1)^j |\underline{q}|^{2j} \partial_{q_0}^{2j+1}A_0(q_0)}{2^{2j}j! \Gamma \left(\frac{3}{2}+j\right)}\\
\nonumber
&&+ \frac{\sqrt{\pi}}{4} \sum_{j=0}^\infty \frac{(-1)^j (|\underline{q}| \partial_{q_0})^{2j+2} A_0(q_0)}{2^{2j}j! \Gamma \left(\frac{5}{2}+j\right)}+\frac{3 \sqrt{\pi}q_0}{2}\sum_{j=0}^\infty \frac{(-1)^j (|\underline{q}| \partial_{q_0})^{2j}A_1(q_0)}{2^{2j}j! \Gamma \left(\frac{3}{2}+j\right)}\\
\label{p2}
&&-\frac{3 \sqrt{\pi}}{4} \sum_{j=0}^\infty \frac{(-1)^j | \underline{q}|^{2j+2} \partial_{q_0}^{2j+1}A_1(q_0)}{2^{2j} j! \Gamma \left( \frac{5}{2}+j\right)},
\end{eqnarray}
and
\begin{equation}
\label{p3}
GCK_0[g_2(q_0)](q)= \frac{\sqrt{\pi}}{2} \sum_{j=0}^\infty \frac{(-1)^j | \underline{q}|^{2j} \partial_{q_0}^{2j+1}A_0(q_0)}{2^{2j}j! \Gamma \left(j+\frac{3}{2}\right)}- \frac{3 \sqrt{\pi}}{2} \sum_{j=0}^\infty \frac{(-1)^j (|\underline{q}| \partial_{q_0})^{2j}A_1(q_0)}{2^{2j}j! \Gamma \left(\frac{3}{2}+j\right)}.
\end{equation}

Thus by \eqref{p2} and \eqref{p3} we have
\begin{eqnarray}
\nonumber
GCK_0[g_1(q_0)](q)+q_0GCK_0[g_2(q_0)](q)&=&\frac{\sqrt{\pi}}{2} \sum_{j=0}^\infty \frac{(-1)^j (|\underline{q}| \partial_{q_0})^{2j}A_0(q_0)}{2^{2j}j! \Gamma \left(\frac{3}{2}+j\right)}\\
\nonumber
&&+ \frac{\sqrt{\pi}}{4} \sum_{j=0}^\infty \frac{(-1)^j (|\underline{q}| \partial_{q_0})^{2j+2} A_0(q_0)}{2^{2j}j! \Gamma \left(\frac{5}{2}+j\right)}\\
\label{scal}
	&&-\frac{3 \sqrt{\pi}}{4} \sum_{j=0}^\infty \frac{(-1)^j | \underline{q}|^{2j+2} \partial_{q_0}^{2j+1}A_1(q_0)}{2^{2j} j! \Gamma \left( \frac{5}{2}+j\right)}.
\end{eqnarray}


We now focus on the $1$-vector part of the GCK extension of the functions $g_1(q_0)$ and $g_2(q_0)$, which we denote by $GCK_1$. By the Leibintz formula, for $j \in \mathbb{N}$, we have
\begin{eqnarray*}
(|\underline{q}| \partial_{q_0})^{2j} \partial_{q_0} g_1(q_0)&=&-q_0 | \underline{q}|^{2j} \partial_{q_0}^{2j+2} A_0(q_0)-2j |\underline{q}|^{2j} \partial_{q_0}^{2j+1}A_0(q_0)+3 | \underline{q}|^{2 j} \partial_{q_0}^{2j}A_1(q_0)\\
&&+3 q_0 | \underline{q}|^{2j} \partial_{q_0}^{2j+1}A_1(q_0)+6j | \underline{q}|^{2j} \partial_{q_0}^{2j} A_1(q_0).
\end{eqnarray*}
Hence, we have
\begin{eqnarray}
	\nonumber
GCK_1[g_1(q_0)](q)&=& \frac{\sqrt{\pi}}{4} \sum_{j=0}^\infty \frac{(-1)^j (|\underline{q}| \partial_{q_0})^{2j} \partial_{q_0} g_1(q_0)}{2^{2j}j! \Gamma \left(j+\frac{5}{2}\right)}\\
\nonumber
&=&- \frac{\sqrt{\pi}q_0}{4} \sum_{j=0}^\infty \frac{(-1)^j | \underline{q}|^{2j} \partial_{q_0}^{2j+2}A_0(q_0)}{2^{2j}j! \Gamma \left(j+\frac{5}{2}\right)}+ \frac{\sqrt{\pi}}{8} \sum_{j=0}^\infty \frac{(-1)^j | \underline{q}|^{2j+2} \partial_{q_0}^{2j+3}A_0(q_0)}{2^{2j} j! \Gamma \left(j+\frac{7}{2}\right)}\\
\nonumber
&&+ \frac{3 \sqrt{\pi}}{4} \sum_{j=0}^\infty \frac{(-1)^j | \underline{q}|^{2j} \partial_{q_0}^{2j}A_1(q_0)}{2^{2j}j! \Gamma \left(j+\frac{5}{2}\right)}+ \frac{3 \sqrt{\pi} q_0}{4} \sum_{j=0}^\infty \frac{(-1)^j| \underline{q}|^{2j} \partial_{q_0}^{2j+1}A_1(q_0)}{2^{2j}j! \Gamma \left(j+\frac{5}{2}\right)}\\
\label{p4}
&&-\frac{3\sqrt{\pi}}{8} \sum_{j=0}^\infty \frac{(-1)^j | \underline{q}|^{2j+2} \partial_{q_0}^{2j+2}A_1(q_0)}{2^{2j}j! \Gamma \left(j+\frac{7}{2}\right)},
\end{eqnarray}
and
\begin{equation}
\label{p5}
GCK_1[g_2(q_0)](q)= \frac{\sqrt{\pi}}{4} \sum_{j=0}^\infty \frac{(-1)^j |\underline{q}|^{2j} \partial_{q_0}^{2j+2}A_0(q_0)}{2^{2j}j!\Gamma \left(j+\frac{5}{2}\right)}- \frac{3 \sqrt{ \pi}}{4} \sum_{j=0}^{\infty} \frac{(-1)^j | \underline{q}|^{2j} \partial_{q_0}^{2j+1}A_1(q_0)}{2^{2j} j! \Gamma \left(j+ \frac{5}{2}\right)}.
\end{equation}

Thus by \eqref{p4} and \eqref{p5} we have
\begin{eqnarray}
\nonumber
GCK_1[g_1(q_0)](q)+q_0 GCK_1[g_2(q_0)](q)&=&\frac{\sqrt{\pi}}{8} \sum_{j=0}^\infty \frac{(-1)^j | \underline{q}|^{2j+2} \partial_{q_0}^{2j+3}A_0(q_0)}{2^{2j} j! \Gamma \left(j+\frac{7}{2}\right)}\\
\nonumber
&&+ \frac{3 \sqrt{\pi}}{4} \sum_{j=0}^\infty \frac{(-1)^j | \underline{q}|^{2j} \partial_{q_0}^{2j}A_1(q_0)}{2^{2j}j! \Gamma \left(j+\frac{5}{2}\right)}\\
\label{vec}
&&-\frac{3\sqrt{\pi}}{8} \sum_{j=0}^\infty \frac{(-1)^j | \underline{q}|^{2j+2} \partial_{q_0}^{2j+2}A_1(q_0)}{2^{2j}j! \Gamma \left(j+\frac{7}{2}\right)}.
\end{eqnarray}
Hence by \eqref{scal} and \eqref{vec} we have
\begin{eqnarray*}
\mathrm{GCK}[g_1(q_0)](q) + q_0\, \mathrm{GCK}[g_2(q_0)](q)&=&GCK_0[g_1(q_0)](q)+q_0GCK_0[g_2(q_0)](q)\\
&&+ \underline{\omega} | \underline{q}| GCK_1[g_1(q_0)](q)+ \underline{\omega}q_0|\underline{q}| GCK_1[g_2(q_0)](q)\\
&=& \alpha(q)+ \underline{\omega} \beta(q),
\end{eqnarray*}
where $\alpha(q)$ and $\beta(q)$ have been introduced in \eqref{real} and in \eqref{imm}, respectively. Finally by \eqref{funct} we get the final result.


\end{proof}

\end{document}
